\chapter{Data Understanding ed Esplorazione dei Dati}
\label{ch:data-understanding}

\section{Obiettivi del Data Understanding}

La fase di comprensione dei dati (\textit{Data Understanding}) costituisce un prerequisito indispensabile per qualsiasi attività di Data Mining. Prima di selezionare o applicare qualsiasi algoritmo, l'analista deve acquisire una visione complessiva della struttura, della distribuzione e della qualità intrinseca del dataset rispetto agli obiettivi del progetto.

In questa fase preliminare è necessario rispondere a domande metodologiche cruciali:
\begin{itemize}
    \item Quali tipologie di attributi compongono il dataset?
    \item Qual è il livello complessivo di qualità dei dati (presenza di rumore, errori o duplicati)?
    \item Come sono distribuite le variabili numeriche e categoriche?
    \item Sono presenti correlazioni evidenti o ridondanze tra gli attributi?
    \item Come sono strutturati gli outlier e quale significato applicativo assumono?
    \item Quali strategie adottare per i valori mancanti?
    \item È necessario derivare o trasformare attributi per rendere più efficace il mining?
\end{itemize}

\section{Tipologie di Dataset}

I dati raccolti si presentano in diverse strutture logico-geometriche in base alla natura del dominio applicativo:

\begin{enumerate}
    \item \textbf{Dati di tipo Record:}
    \begin{itemize}
        \item \textbf{Data Matrix:} Insieme di oggetti descritti dallo stesso insieme fissato di attributi numerici. Un dataset di $m$ oggetti e $n$ attributi è rappresentato formalmente da una matrice $m \times n$:
        \begin{equation}
        \mathbf{X} = \begin{pmatrix}
        x_{1,1} & x_{1,2} & \cdots & x_{1,n} \\
        x_{2,1} & x_{2,2} & \cdots & x_{2,n} \\
        \vdots & \vdots & \ddots & \vdots \\
        x_{m,1} & x_{m,2} & \cdots & x_{m,n}
        \end{pmatrix}
        \end{equation}
        In tale formalizzazione, ciascun oggetto $x_i$ è un punto in uno spazio vettoriale $n$-dimensionale $\mathbb{R}^n$.
        \item \textbf{Document Data (Bag-of-Words):} Ciascun documento testuale è rappresentato da un vettore di termini (\textit{term vector}). Ciascuna componente del vettore corrisponde a una parola del vocabolario, e il valore associato ne indica la frequenza di occorrenza nel testo.
        \item \textbf{Dati Transazionali (\textit{Transaction Data}):} Ciascun record rappresenta una transazione (identificata da un TID) e contiene un insieme variabile di articoli acquistati o eventi osservati.
    \end{itemize}

    \item \textbf{Dati a Grafo (\textit{Graph Data}):} Strutture composte da nodi ed archi che modellano relazioni e connettività. Esempi includono la rete di collegamenti ipertestuali del World Wide Web, i grafi delle reti sociali e le strutture molecolari (ad esempio, la molecola planare del benzene $\text{C}_6\text{H}_6$).

    \item \textbf{Dati Ordinati (\textit{Ordered Data}):}
    \begin{itemize}
        \item \textbf{Sequenze di transazioni:} Successioni temporali di acquisti o eventi per ciascun utente.
        \item \textbf{Dati genomici:} Sequenze biologiche discrete composte dai caratteri dei nucleotidi $\{\text{A}, \text{C}, \text{G}, \text{T}\}$.
        \item \textbf{Dati Spazio-Temporali:} Misurazioni continue o campionate su coordinate geografiche e intervalli temporali regolari (ad esempio, la temperatura media mensile della superficie terrestre e oceanica).
    \end{itemize}
\end{enumerate}

\section{Oggetti e Attributi}

Un dataset è una collezione di \textbf{oggetti dati} (\textit{data objects}) e dei relativi \textbf{attributi} (\textit{attributes}).
\begin{itemize}
    \item Un \textbf{oggetto} (noto anche come record, punto, caso, campione, entità o istanza) rappresenta una singola entità del dominio di interesse.
    \item Un \textbf{attributo} (noto anche come variabile, campo, caratteristica, dimensione o feature) rappresenta una proprietà o caratteristica misurabile di un oggetto (ad esempio, l'età, il reddito imponibile o il colore degli occhi).
\end{itemize}

\subsection{Scale di Misura degli Attributi}
Il tipo di scala di misura di un attributo è definito dall'insieme di proprietà algebriche e operatori matematici che conservano il significato della misura:
\begin{enumerate}
    \item \textbf{Distinzione} ($=, \neq$): Possibilità di distinguere se due valori sono uguali o diversi.
    \item \textbf{Ordine} ($<, >$): Possibilità di stabilire una relazione d'ordine tra valori.
    \item \textbf{Differenze significative} ($+, -$): Presenza di un'unità di misura costante che consente di quantificare l'intervallo tra due valori (lo zero è convenzionale/arbitrario).
    \item \textbf{Rapporti significativi} ($*, /$): Presenza di uno \textbf{zero assoluto intrinseco} (indicante la totale assenza della grandezza misurata), che rende significativo il confronto proporzionale tra grandezze.
\end{enumerate}

\begin{table}[htbp]
\centering
\footnotesize
\begin{tabularx}{\textwidth}{l p{2.2cm} X p{2.6cm} p{3.0cm}}
\toprule
\textbf{Scala} & \textbf{Proprietà} & \textbf{Descrizione} & \textbf{Esempi} & \textbf{Operatori Statistici} \\
\midrule
\textbf{Nominale} & Distinzione ($=, \neq$) & Valori qualitativi o etichette categoriche senza ordinamento intrinseco & Codice fiscale, CAP, colore occhi, genere & Moda, entropia, test $\chi^2$, contingenza \\
\addlinespace
\textbf{Ordinale} & Distinzione e Ordine ($=, \neq, <, >$) & Valori con ordinamento significativo ma intervallo non costante & Titolo di studio, scala Likert (1--10), durezza minerali & Mediana, percentili, correlazione di rango (Spearman) \\
\addlinespace
\textbf{Intervallo} & Ordine e Differenze ($+, -$) & Differenze quantitative dotate di senso; lo zero è convenzionale & Temperatura in $^\circ\text{C}$ o $^\circ\text{F}$, date di calendario & Media, varianza, deviazione standard, Pearson, test $t$ ed $F$ \\
\addlinespace
\textbf{Rapporto} & Tutte le 4 proprietà ($=, <, +, *, /$) & Presenza di zero assoluto; rapporti tra misure fisicamente validi & Età, reddito, massa, conteggi, Kelvin & Media geometrica, armonica, variazione percentuale \\
\bottomrule
\end{tabularx}
\caption{Tassonomia delle scale di misura degli attributi (Stevens).}
\label{tab:attribute-types}
\end{table}

\subsection{Ulteriori Distinzioni degli Attributi}
\begin{itemize}
    \item \textbf{Attributi Binari:} Attributi nominali che ammettono esclusivamente due stati ($0$ e $1$). Si distinguono in:
    \begin{itemize}
        \item \textit{Binari Simmetrici:} Entrambi gli esiti hanno pari importanza informativa (ad esempio il genere sessuale).
        \item \textit{Binari Asimmetrici:} Uno dei due esiti è raro ma di estrema rilevanza informativa (ad esempio, l'esito positivo a un test diagnostico per una malattia rara). Per convenzione, si associa il valore $1$ all'evento di maggiore interesse.
    \end{itemize}
    \item \textbf{Attributi Discreti vs Continui:}
    \begin{itemize}
        \item \textit{Discreti:} Assumono un insieme finito o numerabile di valori (tipicamente rappresentati da numeri interi, conteggi o categorie). Gli attributi binari costituiscono un caso speciale di attributo discreto.
        \item \textit{Continui:} Assumono valori in un sottoinsieme continuo dei numeri reali $\mathbb{R}$ (temperatura, peso, coordinate). Nella pratica computazionale, vengono misurati e memorizzati mediante rappresentazioni in virgola mobile (\textit{floating-point}) con precisione finita.
    \end{itemize}
\end{itemize}

\section{Qualità del Dato (Data Quality)}

La qualità del dato è un fattore critico per il successo di qualunque iniziativa analitica. Come evidenziato negli studi di Thomas Redman, i problemi derivanti da dati di scarsa qualità costano tipicamente alle organizzazioni tra il $10\%$ e il $20\%$ dei loro ricavi complessivi.
Se un modello di rischio creditizio viene addestrato su dati viziati o errati, il sistema rifiuterà debitori solvibili ed erogherà prestiti a individui insolventi.

\subsection{Dimensioni della Qualità del Dato}
I problemi di qualità si manifestano sotto diverse dimensioni:
\begin{enumerate}
    \item \textbf{Accuratezza Sintattica (\textit{Syntactic Accuracy}):} Il valore registrato viola i vincoli formali o non appartiene al dominio prefissato (ad esempio, la stringa \texttt{"fmale"} per il sesso, o la presenza di caratteri testuali in un campo numerico). Tali anomalie sono identificabili in modo immediato ed automatizzato mediante semplici controlli di schema e tipi.
    \item \textbf{Accuratezza Semantica (\textit{Semantic Accuracy}):} Il dato è formalmente corretto e conforme al tipo previsto, ma falso nella realtà di riferimento (ad esempio, associare il nome anagrafico \texttt{"John Smith"} al genere femminile, o registrare un'età biologica pari a $180$ anni). Rilevare queste discrepanze richiede la definizione di regole di dominio (\textit{business rules}) o incroci con sorgenti esterne certificate.
    \item \textbf{Completezza (\textit{Completeness}):} Mancanza sistematica di valori per determinati attributi o assenza di interi record. Una violazione insidiosa della completezza è il \textbf{bias di campionamento} (ad esempio, l'analisi del rischio di insolvenza svolta su un archivio storico che include esclusivamente clienti a cui il credito era già stato concesso, escludendo a priori i profili a basso reddito respinti).
    \item \textbf{Dati Sbilanciati (\textit{Unbalanced Data}):} Distribuzione fortemente asimmetrica delle classi (ad esempio, transazioni fraudolente pari allo $0.1\%$ del totale o prodotti difettosi su linee di produzione ad altissima automazione).
    \item \textbf{Tempestività (\textit{Timeliness}):} Obsoleto grado di aggiornamento dei dati rispetto alla dinamica corrente del fenomeno osservato.
    \item \textbf{Dati Duplicati (\textit{Duplicate Data}):} Presenza di record identici o quasi-identici relativi alla medesima entità (ad esempio, un medesimo cliente registrato più volte con indirizzi email differenti a seguito della fusione di due banche dati eterogenee). Tali problematiche richiedono processi dedicati di deduplicazione e \textit{record linkage}.
\end{enumerate}

\begin{noteblock}{Anomalie nei Flussi di Misurazione e Sensori}
Nei sistemi di monitoraggio industriale e IoT, un sensore guasto o bloccato trasmette tipicamente un valore costante pari a zero. In fase di esplorazione, tali registrazioni non devono essere considerate misurazioni fisiche effettive, bensì interpretate e gestite a tutti gli effetti come \textbf{valori mancanti} (dovuti a malfunzionamento del dispositivo).
\end{noteblock}

\section{Statistica Descrittiva ed Esplorazione Univariata}

\subsection{Misure di Tendenza Centrale}
\begin{itemize}
    \item \textbf{Media Campionaria ($\bar{x}$):} Per un campione di dimensione $m$:
    \begin{equation}
    \bar{x} = \frac{1}{m} \sum_{i=1}^m x_i
    \end{equation}
    La media è una misura distributiva (può essere calcolata partizionando il dataset in sottoinsiemi), ma è \textbf{estremamente sensibile alla presenza di outlier}.
    \item \textbf{Mediana:} Il valore che occupa la posizione centrale nell'ordinamento crescente delle osservazioni (o la media dei due valori centrali se $m$ è pari). È una misura robusta, con un punto di rottura (\textit{breakdown point}) del $50\%$.
    \item \textbf{Media Troncata (\textit{Trimmed Mean}):} Media calcolata escludendo una percentuale fissa (ad esempio il $5\%$ o il $10\%$) dei valori più estremi a entrambe le estremità.
    \item \textbf{Moda:} Il valore o la categoria che si presenta con la frequenza assoluta massima. Una distribuzione può essere unimodale, bimodale, multimodale, oppure priva di moda se tutte le osservazioni compaiono una sola volta.
\end{itemize}

\subsection{Asimmetria e Multimodalità}
\begin{itemize}
    \item \textbf{Asimmetria Positiva (\textit{Positively Skewed / Right-skewed}):} La coda destra della distribuzione è allungata. La maggior parte delle osservazioni è concentrata su valori inferiori alla media aritmetica ($\text{Moda} < \text{Mediana} < \text{Media}$). Un esempio tipico è la distribuzione dei \textbf{redditi}: la grande maggioranza dei cittadini percepisce redditi bassi o medi, mentre un esiguo numero di individui percepisce redditi estremamente elevati, trainando la media verso l'alto.
    \item \textbf{Distribuzione Bimodale:} Presenza di due picchi distinti nella densità di frequenza. Segnala tipicamente l'influenza di una variabile latente binaria non esplicitata (ad esempio, la distribuzione del peso corporeo umano in una popolazione mista, influenzata dalla variabile binaria sesso).
\end{itemize}

\subsection{Misure di Dispersione}
\begin{itemize}
    \item \textbf{Range:} Differenza tra valore massimo e valore minimo: $\text{Range} = x_{\max} - x_{\min}$.
    \item \textbf{Varianza Campionaria ($s^2$) e Deviazione Standard ($s$):}
    \begin{equation}
    s^2 = \frac{1}{m-1} \sum_{i=1}^m (x_i - \bar{x})^2, \quad s = \sqrt{s^2}
    \end{equation}
    La deviazione standard misura la dispersione attorno alla media; è nulla solo se tutte le osservazioni sono identiche ($s=0$), altrimenti $s > 0$. Risulta anch'essa sensibile agli outlier.
    \item \textbf{Misure Robuste di Dispersione:}
    \begin{itemize}
        \item \textit{Absolute Average Deviation} (AAD):
        \begin{equation}
        \text{AAD} = \frac{1}{m} \sum_{i=1}^m |x_i - \bar{x}|
        \end{equation}
        \item \textit{Median Absolute Deviation} (MAD):
        \begin{equation}
        \text{MAD} = \operatorname{mediana}\big(|x_i - \operatorname{mediana}(x)|\big)
        \end{equation}
        \item \textit{Interquartile Range} (IQR): Distanza tra il terzo quartile $Q_3$ ($75^\circ$ percentile) e il primo quartile $Q_1$ ($25^\circ$ percentile):
        \begin{equation}
        \text{IQR} = Q_3 - Q_1
        \end{equation}
    \end{itemize}
\end{itemize}

\subsection{Box Plot e Sintesi a Cinque Numeri}
Il \textbf{Box Plot} (grafico a scatola e baffi) offre una sintesi visiva non parametrica della distribuzione basata sulla \textbf{sintesi a cinque numeri}:
\[
\big[\text{Minimo}, \; Q_1, \; \text{Mediana} \, (Q_2), \; Q_3, \; \text{Massimo}\big]
\]
La costruzione geometrica si fonda sul concetto di quantile: dato $p \in (0, 100)$, il \textbf{$p\%$-quantile} (o $p$-esimo percentile) \`e la soglia $x_p$ tale per cui il $p\%$ delle osservazioni campionarie assume valori inferiori o uguali a $x_p$, mentre il $(100-p)\%$ assume valori superiori. In particolare:
\begin{itemize}
    \item $Q_1$ corrisponde al $25\%$-quantile (primo quartile).
    \item La mediana $Q_2$ corrisponde al $50\%$-quantile (secondo quartile).
    \item $Q_3$ corrisponde al $75\%$-quantile (terzo quartile).
    \item La larghezza della scatola centrale riflette l'intervallo interquartile $\text{IQR} = Q_3 - Q_1$, racchiudendo il $50\%$ centrale della massa dei dati.
\end{itemize}

\begin{figure}[htbp]
\centering
\begin{tikzpicture}[xscale=0.9, yscale=0.85]
    % Axis
    \draw[thick, -Stealth] (0,0) -- (12,0) node[right, font=\small] {Valori};
    \foreach \x in {1,3,5,7,9,11} \draw (\x, 0.1) -- (\x, -0.1) node[below, font=\scriptsize] {\x};

    % Box (Q1=4, Mediana=6, Q3=8)
    \draw[thick, fill=blue!15] (4,1) rectangle (8,3);
    \draw[very thick, red!80!black] (6,1) -- (6,3) node[above, font=\scriptsize\bfseries] {Mediana ($Q_2$)};

    % Whiskers
    \draw[thick] (4,2) -- (2,2);
    \draw[thick] (2,1.5) -- (2,2.5) node[above, font=\scriptsize] {Baffo Min};
    \draw[thick] (8,2) -- (10,2);
    \draw[thick] (10,1.5) -- (10,2.5) node[above, font=\scriptsize] {Baffo Max};

    % Outliers
    \fill[red] (11.2, 2) circle (3pt) node[above, font=\scriptsize] {Outlier};
    \fill[red] (0.8, 2) circle (3pt) node[above, font=\scriptsize] {Outlier};

    % Annotations
    \draw[Stealth-Stealth, thick] (4,0.5) -- node[below, font=\scriptsize] {$\text{IQR} = Q_3 - Q_1$} (8,0.5);
    \node[font=\scriptsize] at (4,3.3) {$Q_1$ (25\%)};
    \node[font=\scriptsize] at (8,3.3) {$Q_3$ (75\%)};
\end{tikzpicture}
\caption{Anatomia geometrica di un Box Plot con individuazione degli outlier.}
\label{fig:boxplot-anatomy}
\end{figure}

I limiti dei baffi sono calcolati convenzionalmente estendendosi fino al valore osservato pi\`u estremo entro l'intervallo $[Q_1 - 1.5 \times \text{IQR}, \; Q_3 + 1.5 \times \text{IQR}]$. Le osservazioni che cadono al di fuori di tali soglie sono etichettate visivamente come \textbf{outlier}.

Un'estensione fondamentale per l'analisi bivariata e multivariata \`e costituita dai \textbf{Box Plot Condizionali} (\textit{Conditional Box Plots}): la distribuzione di un attributo numerico continuo viene stratificata e visualizzata affiancando i box plot per ciascuna modalit\`a o categoria di una variabile discreta (tipicamente la classe target o un raggruppamento naturale). In Figura~\ref{fig:conditional-boxplot} \`e illustrata la distribuzione della lunghezza del petalo (\textit{Petal Length}) stratificata sulle tre specie del dataset Iris: il grafico evidenzia visivamente che la specie \textit{Setosa} \`e nettamente separabile dalle restanti due con variabilit\`a molto ridotta, mentre \textit{Versicolor} e \textit{Virginica} presentano una parziale sovrapposizione nei rispettivi intervalli interquartili.

\begin{figure}[htbp]
\centering
\includegraphics[width=0.88\textwidth]{iris_conditional_boxplot.pdf}
\caption{Box Plot Condizionale calcolato sui dati reali del dataset Iris: distribuzione empirica della lunghezza del petalo (\textit{Petal Length}) stratificata sulle tre specie, con evidenza della mediana, dei quartili e dei singoli campioni osservati.}
\label{fig:conditional-boxplot}
\end{figure}

\begin{alertblock}{Confronto Critico: Box Plot vs Istogramma}
Sebbene il box plot sintetizzi efficacemente posizione e dispersione, \textbf{non \`e in grado di mostrare la reale forma della distribuzione}. Due distribuzioni aventi densit\`a radicalmente differenti (ad esempio, una distribuzione unimodale uniforme e una distribuzione fortemente bimodale con una valle centrale vuota) possono produrre identici valori di $Q_1$, mediana, $Q_3$, Minimo e Massimo, e generare pertanto lo stesso identico box plot. L'istogramma \`e indispensabile per rilevare la multimodalit\`a.
\end{alertblock}

\begin{figure}[htbp]
\centering
\begin{minipage}{0.46\textwidth}
\centering
\begin{tikzpicture}[scale=0.65]
    \draw[thick, -Stealth] (0,0) -- (5.5,0) node[right, font=\tiny] {$x$};
    \draw[thick, -Stealth] (0,0) -- (0,3.5) node[above, font=\tiny] {Frequenza};
    % Uniform / unimodal shape
    \draw[fill=orange!30, draw=orange!80, thick] (0.5,0) rectangle (1.3, 1.0);
    \draw[fill=orange!30, draw=orange!80, thick] (1.3,0) rectangle (2.1, 2.2);
    \draw[fill=orange!30, draw=orange!80, thick] (2.1,0) rectangle (2.9, 3.0);
    \draw[fill=orange!30, draw=orange!80, thick] (2.9,0) rectangle (3.7, 2.2);
    \draw[fill=orange!30, draw=orange!80, thick] (3.7,0) rectangle (4.5, 1.0);
    \node[font=\scriptsize\bfseries] at (2.5, -0.6) {Densit\`a Unimodale};
\end{tikzpicture}
\end{minipage}\hfill
\begin{minipage}{0.46\textwidth}
\centering
\begin{tikzpicture}[scale=0.65]
    \draw[thick, -Stealth] (0,0) -- (5.5,0) node[right, font=\tiny] {$x$};
    \draw[thick, -Stealth] (0,0) -- (0,3.5) node[above, font=\tiny] {Frequenza};
    % Bimodal shape with empty center
    \draw[fill=purple!30, draw=purple!80, thick] (0.5,0) rectangle (1.3, 2.8);
    \draw[fill=purple!30, draw=purple!80, thick] (1.3,0) rectangle (2.1, 1.2);
    \draw[fill=purple!30, draw=purple!80, thick] (2.1,0) rectangle (2.9, 0.3);
    \draw[fill=purple!30, draw=purple!80, thick] (2.9,0) rectangle (3.7, 1.2);
    \draw[fill=purple!30, draw=purple!80, thick] (3.7,0) rectangle (4.5, 2.8);
    \node[font=\scriptsize\bfseries] at (2.5, -0.6) {Densit\`a Bimodale};
\end{tikzpicture}
\end{minipage}

\vspace{0.3cm}
\begin{tikzpicture}[scale=0.75]
    \draw[thick, -Stealth] (0,0) -- (8.5,0) node[right, font=\scriptsize] {$x$};
    \foreach \x/\val in {0.8/Min, 2.4/$Q_1$, 4.0/Mediana, 5.6/$Q_3$, 7.2/Max} {
        \draw (\x, 0.1) -- (\x, -0.1) node[below, font=\tiny] {\val};
    }
    \draw[thick, fill=gray!25] (2.4,0.4) rectangle (5.6,1.2);
    \draw[very thick, red!80!black] (4.0,0.4) -- (4.0,1.2);
    \draw[thick] (2.4,0.8) -- (0.8,0.8); \draw[thick] (0.8,0.5) -- (0.8,1.1);
    \draw[thick] (5.6,0.8) -- (7.2,0.8); \draw[thick] (7.2,0.5) -- (7.2,1.1);
    \node[font=\scriptsize, red!80!black] at (4.0, 1.55) {Identico Box Plot per entrambe le distribuzioni!};
\end{tikzpicture}
\caption{Confronto critico: due distribuzioni con andamenti diametralmente opposti (unimodale vs bimodale con centro svuotato) generano la medesima sintesi a cinque numeri e lo stesso Box Plot.}
\label{fig:boxplot-vs-hist}
\end{figure}

\subsection{Bar Chart per Attributi Categorici}
Il \textbf{Bar Chart} (grafico a barre) rappresenta visivamente la distribuzione di frequenza di un attributo qualitativo o categorico. Ciascuna categoria discreta \`e rappresentata da una barra verticale (o orizzontale), la cui altezza corrisponde al conteggio assoluto o alla frequenza relativa. A differenza degli istogrammi, le barre sono separate da spazi vuoti per evidenziare che le modalit\`a non costituiscono un continuo numerico.

\begin{figure}[htbp]
\centering
\begin{tikzpicture}[xscale=1.4, yscale=0.07]
    \draw[thick, -Stealth] (0,0) -- (4.5,0) node[right, font=\small] {Stato Civile};
    \draw[thick, -Stealth] (0,0) -- (0,55) node[above, font=\small] {Frequenza};
    \foreach \y in {10, 20, 30, 40, 50} \draw (0.08, \y) -- (-0.08, \y) node[left, font=\scriptsize] {\y};

    \draw[fill=blue!35, draw=blue!80, thick] (0.6,0) rectangle (1.4,42);
    \node[above, font=\scriptsize\bfseries] at (1.0, 42) {42};
    \node[below, font=\small] at (1.0, -2) {Single};

    \draw[fill=blue!35, draw=blue!80, thick] (1.8,0) rectangle (2.6,33);
    \node[above, font=\scriptsize\bfseries] at (2.2, 33) {33};
    \node[below, font=\small] at (2.2, -2) {Married};

    \draw[fill=blue!35, draw=blue!80, thick] (3.0,0) rectangle (3.8,18);
    \node[above, font=\scriptsize\bfseries] at (3.4, 18) {18};
    \node[below, font=\small] at (3.4, -2) {Divorced};
\end{tikzpicture}
\caption{Esempio di Bar Chart per un attributo categorico nominale (Stato Civile).}
\label{fig:barchart-example}
\end{figure}

\subsection{Istogrammi per Attributi Numerici e Scelta dei Bin}
Un \textbf{istogramma} mostra la distribuzione empirica di frequenza di un attributo numerico continuo. Il dominio di variabilit\`a della grandezza viene partizionato (\textit{discretizzato}) in un numero prefissato $k$ di intervalli contigui detti \textit{bin} (o classi di ampiezza). Per ciascun bin, l'altezza della barra sovrastante \`e proporzionale al numero di osservazioni campionarie che cadono al suo interno.

\begin{figure}[htbp]
\centering
\begin{tikzpicture}[xscale=0.9, yscale=0.08]
    \draw[thick, -Stealth] (0,0) -- (11.5,0) node[right, font=\small] {Et\`a};
    \draw[thick, -Stealth] (0,0) -- (0,50) node[above, font=\small] {Frequenza};
    \foreach \y in {10, 20, 30, 40, 50} \draw (0.1, \y) -- (-0.1, \y) node[left, font=\scriptsize] {\y};

    % 6 contiguous bins: 20-30, 30-40, 40-50, 50-60, 60-70, 70-80
    \draw[fill=teal!30, draw=teal!80, thick] (1,0) rectangle (2.5, 14);
    \draw[fill=teal!30, draw=teal!80, thick] (2.5,0) rectangle (4, 44);
    \draw[fill=teal!30, draw=teal!80, thick] (4,0) rectangle (5.5, 32);
    \draw[fill=teal!30, draw=teal!80, thick] (5.5,0) rectangle (7, 20);
    \draw[fill=teal!30, draw=teal!80, thick] (7,0) rectangle (8.5, 11);
    \draw[fill=teal!30, draw=teal!80, thick] (8.5,0) rectangle (10, 5);

    \foreach \x/\label in {1/20, 2.5/30, 4/40, 5.5/50, 7/60, 8.5/70, 10/80} {
        \draw (\x, 1.5) -- (\x, -1.5) node[below, font=\scriptsize] {\label};
    }
\end{tikzpicture}
\caption{Esempio di Istogramma per un attributo numerico continuo con bin contigui di ampiezza costante.}
\label{fig:histogram-example}
\end{figure}

Una scelta scorretta del numero di bin altera profondamente la percezione dei dati:
\begin{itemize}
    \item Un numero insufficiente di bin provoca \textit{oversmoothing}, livellando la distribuzione e mascherando asimmetrie e bimodalit\`a.
    \item Un numero eccessivo di bin provoca \textit{undersmoothing}, frammentando il campione in frequenze discontinue ed enfatizzando il rumore campionario.
\end{itemize}

Per distribuzioni unimodali approssimativamente normali e campioni di dimensione moderata $n$, la stima standard del numero ottimale di intervalli \`e fornita dalla \textbf{regola di Sturges}:
\begin{equation}
k = 1 + \lceil \log_2(n) \rceil = 1 + \lceil 3.322 \log_{10}(n) \rceil
\end{equation}
dove $n$ rappresenta la cardinalit\`a campionaria. Tale euristica assume esplicitamente una distribuzione di tipo gaussiano. Quando il volume dei dati cresce notevolmente o la distribuzione sottostante presenta pi\`u modi o forte asimmetria, la regola di Sturges tende a generare un numero insufficiente di classi; in tali casi risulta fondamentale testare sperimentalmente diversi valori di $k$ (ad esempio confrontando partizioni a bassa, media ed elevata risoluzione, come 5, 17 o 50 bin) per verificare la stabilit\`a dei pattern riscontrati.

\section{Tecniche di Visualizzazione Multivariata}

\subsection{Scatter Plot e Visualizzazione a Più Dimensioni}
Lo \textbf{Scatter Plot} (grafico a dispersione) tratta coppie di attributi numerici come coordinate cartesiane $(x_i, y_i)$ nel piano. \`E il principale strumento visuale per:
\begin{itemize}
    \item Identificare correlazioni lineari o andamenti funzionali non lineari.
    \item Rilevare addensamenti naturali di punti (\textit{cluster}).
    \item Individuare outlier multivariati non visibili nelle proiezioni univariate.
\end{itemize}

\begin{example}[Analisi del Dataset Iris di Anderson/Fisher]
Il celebre dataset \textit{Iris} (150 campioni di fiori appartenenti a tre specie: \textit{Setosa}, \textit{Versicolor}, \textit{Virginica}) descrive 4 attributi numerici: lunghezza e larghezza del sepalo, lunghezza e larghezza del petalo.
Lo scatter plot bivariato mostra che gli attributi relativi al petalo (\textit{Petal Length} vs \textit{Petal Width}) garantiscono una separazione netta tra le specie, consentendo inoltre di isolare eventuali anomalie.
\end{example}

\begin{figure}[htbp]
\centering
\includegraphics[width=0.82\textwidth]{iris_scatterplot_real.pdf}
\caption{Scatter Plot bivariato reale (\textit{Petal Length} vs \textit{Petal Width}) per tutti i 150 campioni del dataset Iris, con codifica cromatica e geometrica per classe e identificazione visiva della netta separazione di \textit{Setosa} e dei punti atipici di confine.}
\label{fig:scatterplot-outlier}
\end{figure}

Per esplorare simultaneamente tutte le coppie di attributi numerici, si impiega una \textbf{Scatter Plot Matrix} (SPLOM), ovvero una griglia $n \times n$ in cui ogni cella non diagonale $(i, j)$ visualizza il grafico di dispersione tra l'attributo $i$ e l'attributo $j$, mentre la diagonale principale ospita gli istogrammi o le stime di densit\`a univariate.

\begin{figure}[htbp]
\centering
\includegraphics[width=0.92\textwidth]{iris_splom_real.pdf}
\caption{Scatter Plot Matrix (SPLOM) calcolata su tutti i 150 campioni del dataset Iris attraverso i 4 attributi continui, con stima di densit\`a univariata (KDE) sulla diagonale principale e grafici di dispersione bivariati nelle celle esterne.}
\label{fig:splom-example}
\end{figure}

\subsection{La Visualizzazione come Test ed Euristica Esplorativa}
L'ispezione visiva rappresenta un primo insostituibile test esplorativo sulla presenza di struttura nei dati, la cui corretta interpretazione richiede una precisa consapevolezza metodologica:
\begin{itemize}
    \item \textbf{Presenza di evidenza grafica:} Quando le visualizzazioni evidenziano pattern netti, raggruppamenti separati (\textit{cluster}) o eccezioni isolate (\textit{outlier}), si ha la certezza empirica che vi sia una propriet\`a strutturale reale o un fenomeno oggettivo all'interno del dataset.
    \item \textbf{Assenza di evidenza grafica non implica assenza di struttura:} Al contrario, quando le visualizzazioni ordinarie non mostrano alcuna configurazione apparente o sembrano puramente casuali, \textbf{non} si pu\`o dedurre che i dati siano privi di relazioni. Relazioni multivariate complesse, deboli associazioni non lineari o dipendenze spaziali che coinvolgono quattro o pi\`u attributi contemporaneamente sfuggono intrinsecamente a proiezioni bidimensionali, richiedendo l'impiego di tecniche analitiche o proiezioni a dimensionalit\`a superiore.
\end{itemize}

\subsection{Visualizzazione ad Alta Dimensionalità}

\subsubsection{Coordinate Parallele (Parallel Coordinates)}
Quando il numero di dimensioni supera 3, la rappresentazione cartesiana convenzionale risulta inapplicabile. Le \textbf{Coordinate Parallele} sostituiscono gli assi ortogonali con un insieme di assi paralleli verticali equidistanti, ciascuno corrispondente a un attributo.

\begin{figure}[htbp]
\centering
\includegraphics[width=0.88\textwidth]{iris_parallel_coordinates.pdf}
\caption{Rappresentazione a 4 dimensioni mediante Coordinate Parallele calcolata sui dati reali del dataset Iris (150 campioni normalizzati in $[0, 1]$): le curve sottili rappresentano le singole osservazioni, mentre le linee marcate evidenziano i profili medi distintivi di ciascuna specie.}
\label{fig:parallel-coordinates}
\end{figure}

\begin{itemize}
    \item Ogni oggetto multidimensionale \`e rappresentato da una \textbf{linea spezzata poligonale} che connette i rispettivi valori scalati lungo gli assi verticali.
    \item Gli oggetti appartenenti alla stessa classe tendono a formare fasci di linee parallele o con andamenti coordinati.
    \item \textbf{Sensibilità all'ordinamento degli assi:} La visibilit\`a di raggruppamenti e correlazioni dipende criticamente dalla disposizione sequenziale degli assi. Due attributi fortemente correlati mostrano pattern geometrici regolari solo se posizionati come assi adiacenti.
\end{itemize}

\subsubsection{Radar Plot e Star Plot}
Nei \textbf{Radar Plot}, gli assi partono da un centro comune e si irradiano a stella a intervalli angolari costanti. Il valore di ciascun attributo definisce un vertice lungo il proprio raggio; congiungendo i vertici si ottiene un poligono per ogni istanza.
Gli \textbf{Star Plot} applicano il medesimo principio disegnando ciascun oggetto in un grafico separato, facilitando il confronto per profili individuali.

\begin{figure}[htbp]
\centering
\includegraphics[width=0.62\textwidth]{iris_radar_plot.pdf}
\caption{Radar Plot empirico multivariato calcolato sui profili medi normalizzati delle tre specie di Iris sui 4 attributi morfologici.}
\label{fig:radar-plot-example}
\end{figure}

\section{Analisi di Correlazione e Associazione}

L'analisi di correlazione ha lo scopo di identificare, quantificare e interpretare la forza e la direzione della relazione lineare tra coppie di attributi numerici continui.

\subsection{Motivazioni e Multicollinearit\`a nei Modelli Predittivi}
Nel processo di data understanding, valutare la correlazione tra variabili persegue diversi obiettivi metodologici:
\begin{itemize}
    \item \textbf{Individuazione di ridondanze informative:} Se due attributi manifestano un andamento pressoch\'e coincidente, uno di essi pu\`o essere scartato senza perdita apprezzabile di potere descrittivo, riducendo la dimensionalit\`a dello spazio di lavoro.
    \item \textbf{Preparazione per modelli predittivi e clustering:} Molteplici algoritmi di machine learning (quali regressione lineare, logistic regression o clustering basato su distanze) assumono feature indipendenti o beneficiano fortemente della rimozione di attributi collineari.
    \item \textbf{Rilevamento della multicollinearit\`a:} Si verifica quando due o pi\`u variabili predittive (\textit{feature}) in un modello di regressione risultano fortemente legate da una relazione lineare. Si consideri l'esempio paradigmatico della stima del valore degli immobili impiegando congiuntamente come predittori sia la \textit{superficie in metri quadri} ($m^2$) sia la \textit{superficie in piedi quadri} ($\text{ft}^2$). Le due variabili veicolano esattamente la medesima informazione; di conseguenza, il modello di regressione entra in severa instabilit\`a numerica, non riuscendo ad isolare il contributo marginale di ciascun predittore e gonfiando la varianza degli stimatori.
\end{itemize}

\subsection{Dalla Varianza alla Covarianza Campionaria}
Mentre la varianza campionaria quantifica la dispersione delle osservazioni di una singola variabile attorno alla propria media:
\begin{equation}
s_x^2 = \operatorname{Var}(X) = \frac{1}{n-1} \sum_{i=1}^n (x_i - \bar{x})^2
\end{equation}
nelle applicazioni analitiche l'interesse \`e costantemente rivolto alle interazioni tra coppie di variabili:
\begin{itemize}
    \item I clienti con reddito imponibile pi\`u elevato tendono a spendere somme maggiori?
    \item Il prezzo unitario di un bene determina una contrazione delle quantit\`a vendute?
    \item Un maggior numero di ore di studio settimanali conduce sistematicamente a voti d'esame superiori?
    \item Due sensori fisici disposti in parallelo forniscono misure ampiamente ridondanti?
\end{itemize}

Per misurare come due grandezze numeriche variano simultaneamente, si introduce la \textbf{covarianza campionaria}:

\begin{definition}[Covarianza Campionaria]
Dato un campione di $n$ coppie di valori $(x_i, y_i)$, la covarianza campionaria tra $X$ e $Y$ \`e definita da:
\begin{equation}
\operatorname{Cov}(X, Y) = \frac{1}{n-1} \sum_{i=1}^n (x_i - \bar{x})(y_i - \bar{y})
\end{equation}
dove $\bar{x} = \frac{1}{n}\sum_{i=1}^n x_i$ e $\bar{y} = \frac{1}{n}\sum_{i=1}^n y_i$ rappresentano le medie campionarie.
\end{definition}

Una propriet\`a teorica fondamentale sancisce che la varianza \`e semplicemente la covarianza di una variabile con se stessa:
\begin{equation}
\operatorname{Cov}(X, X) = \frac{1}{n-1} \sum_{i=1}^n (x_i - \bar{x})^2 = \operatorname{Var}(X) = s_x^2
\end{equation}

\subsection{Origine del Segno e Interpretazione Geometrica per Quadranti}
Il segno algebrico di ciascun termine della sommatoria dipende dal prodotto degli scarti congiunti rispetto alle medie:
\begin{itemize}
    \item \textbf{Scarti concordi positivi:} $(x_i - \bar{x}) > 0$ e $(y_i - \bar{y}) > 0 \implies (+)\times(+) = \mathbf{+}$.
    \item \textbf{Scarti concordi negativi:} $(x_i - \bar{x}) < 0$ e $(y_i - \bar{y}) < 0 \implies (-)\times(-) = \mathbf{+}$.
    \item \textbf{Scarti discordi:} $(+)\times(-) = \mathbf{-}$ oppure $(-)\times(+) = \mathbf{-}$.
\end{itemize}

Traslando l'origine cartesiana nel baricentro $(\bar{x}, \bar{y})$, lo spazio viene suddiviso in quattro quadranti (Figura~\ref{fig:covariance-quadrants}):
\begin{enumerate}
    \item Se dominano le osservazioni nei quadranti \textbf{I e III} (scarti concordi), i prodotti positivi prevalgono e la covarianza campionaria \`e positiva ($\operatorname{Cov}(X, Y) > 0$).
    \item Se prevalgono i quadranti \textbf{II e IV} (scarti discordi), dominano i prodotti negativi e la covarianza risulta negativa ($\operatorname{Cov}(X, Y) < 0$).
    \item Se i punti sono dispersi in modo uniforme o simmetrico tra i quattro quadranti, i termini positivi e negativi si elidono reciprocamente, generando una covarianza prossima a zero ($\operatorname{Cov}(X, Y) \approx 0$).
\end{enumerate}

\begin{figure}[htbp]
\centering
\begin{tikzpicture}[scale=0.95]
    % Central cross at mean (x_bar, y_bar)
    \def\xb{3.2}
    \def\yb{2.8}

    % Shaded quadrants
    \fill[green!8] (\xb, \yb) rectangle (6.5, 5.3);
    \fill[green!8] (0.5, 0.5) rectangle (\xb, \yb);
    \fill[red!8] (0.5, \yb) rectangle (\xb, 5.3);
    \fill[red!8] (\xb, 0.5) rectangle (6.5, \yb);

    % Axes
    \draw[thick, -Stealth] (0.5, 0.5) -- (6.8, 0.5) node[right, font=\small] {$X$};
    \draw[thick, -Stealth] (0.5, 0.5) -- (0.5, 5.6) node[above, font=\small] {$Y$};

    % Mean dashed lines
    \draw[dashed, thick, gray!70] (\xb, 0.5) -- (\xb, 5.3);
    \draw[dashed, thick, gray!70] (0.5, \yb) -- (6.5, \yb);
    \node[below, font=\small\bfseries] at (\xb, 0.5) {$\bar{x}$};
    \node[left, font=\small\bfseries] at (0.5, \yb) {$\bar{y}$};

    % Quadrant Labels and signs
    \node[align=center, font=\scriptsize] at (5.0, 4.6) {\textbf{Quadrante I}\\$(x_i > \bar{x}, \; y_i > \bar{y})$\\$(+) \times (+) = \mathbf{+}$};
    \node[align=center, font=\scriptsize] at (1.8, 4.6) {\textbf{Quadrante II}\\$(x_i < \bar{x}, \; y_i > \bar{y})$\\$(-) \times (+) = \mathbf{-}$};
    \node[align=center, font=\scriptsize] at (1.8, 1.4) {\textbf{Quadrante III}\\$(x_i < \bar{x}, \; y_i < \bar{y})$\\$(-) \times (-) = \mathbf{+}$};
    \node[align=center, font=\scriptsize] at (5.0, 1.4) {\textbf{Quadrante IV}\\$(x_i > \bar{x}, \; y_i < \bar{y})$\\$(+) \times (-) = \mathbf{-}$};

    % Sample points illustrating positive covariance
    \foreach \p in {(1.2, 1.1), (1.6, 1.8), (2.1, 1.5), (2.5, 2.2), (2.8, 2.4),
                    (3.6, 3.2), (4.0, 3.8), (4.5, 3.5), (5.1, 4.2), (5.8, 4.7), (6.1, 4.9)} {
        \fill[blue!80!black] \p circle (2.5pt);
    }
    \fill[blue!80!black] (2.2, 3.3) circle (2.5pt);
    \fill[blue!80!black] (4.2, 2.0) circle (2.5pt);

    \node[font=\footnotesize\itshape, blue!80!black] at (3.5, -0.2) {Prevalenza nei Quadranti I e III $\implies \operatorname{Cov}(X, Y) > 0$};
\end{tikzpicture}
\caption{Interpretazione geometrica della covarianza per quadranti centrati in $(\bar{x}, \bar{y})$. I punti nei quadranti concordi I e III producono contributi positivi, mentre i quadranti discordi II e IV producono contributi negativi.}
\label{fig:covariance-quadrants}
\end{figure}

\subsection{Limite della Covarianza: Dipendenza dalla Scala Metrica}
Nonostante la covarianza indichi chiaramente il segno della concordanza lineare, essa presenta un limite critico: \textbf{dipende strettamente dalle unit\`a di misura delle variabili}.
Si consideri ad esempio la relazione tra reddito annuo e spesa:
\begin{itemize}
    \item Se entrambe le variabili sono espresse in euro, l'unit\`a di misura della covarianza \`e pari a $\text{euro}^2$.
    \item Se si converte il reddito in \textit{migliaia di euro}, il valore numerico della covarianza subisce un ridimensionamento di un fattore $1000$, sebbene la relazione sostanziale sottostante tra le grandezze rimanga rigorosamente immutata.
\end{itemize}
Di conseguenza, la covarianza descrive la direzione della dipendenza lineare, ma la sua magnitudo assoluta non \`e comparabile tra coppie di variabili eterogenee n\'e quantifica l'intensit\`a intrinseca del legame.

\subsection{Coefficiente di Correlazione Lineare di Pearson}
Per ottenere una misura standardizzata e priva di unit\`a di misura, si divide la covarianza campionaria per il prodotto delle rispettive deviazioni standard:

\begin{definition}[Coefficiente di Correlazione di Pearson]
Dato un campione di $n$ coppie $(x_i, y_i)$, il coefficiente campionario di correlazione lineare $r_{xy}$ \`e definito come:
\begin{equation}
r_{xy} = \frac{\operatorname{Cov}(X, Y)}{s_x s_y} = \frac{\sum_{i=1}^n (x_i - \bar{x})(y_i - \bar{y})}{\sqrt{\sum_{i=1}^n (x_i - \bar{x})^2 \sum_{i=1}^n (y_i - \bar{y})^2}}
\end{equation}
dove $s_x, s_y$ sono le deviazioni standard campionarie di $X$ e $Y$.
\end{definition}

Propriet\`a cardine del coefficiente di Pearson:
\begin{itemize}
    \item \textbf{Adimensionalit\`a e invarianza di scala:} $r_{xy}$ \`e un numero puro, invariante rispetto a qualunque riscalamento lineare affine positivo ($X' = aX + b$, con $a > 0$).
    \item \textbf{Delimitazione:} Risulta sempre $-1 \le r_{xy} \le +1$.
    \item \textbf{Perfetta collinearit\`a ($|r_{xy}| = 1$):} I punti giacciono esattamente su una linea retta; con pendenza positiva se $r_{xy} = +1$, negativa se $r_{xy} = -1$.
    \item \textbf{Assenza di relazione lineare ($r_{xy} \approx 0$):} Segnala l'assenza di un legame lineare.
    \item \textbf{Cecit\`a alle relazioni non lineari:} Come illustrato in Figura~\ref{fig:pearson-cases}(d), un coefficiente nullo ($r_{xy} \approx 0$) non dimostra affatto indipendenza: relazioni deterministiche paraboliche o quadratiche possono avere correlazione lineare esattamente pari a zero.
\end{itemize}

\begin{figure}[htbp]
\centering
\begin{minipage}{0.47\textwidth}
\centering
\begin{tikzpicture}[scale=0.55]
    \draw[thick, -Stealth] (0,0) -- (5,0) node[right, font=\tiny] {$X$};
    \draw[thick, -Stealth] (0,0) -- (0,5) node[above, font=\tiny] {$Y$};
    \foreach \p in {(0.5,0.7), (1,1.2), (1.5,1.4), (2,2.2), (2.5,2.6), (3,3.1), (3.5,3.7), (4,4.1), (4.5,4.6)} {
        \fill[blue!70] \p circle (2.5pt);
    }
    \draw[red, thick, dashed] (0.3,0.4) -- (4.7,4.8);
    \node[font=\scriptsize\bfseries] at (2.5,-0.8) {(a) $r_{xy} \approx +0.98$};
\end{tikzpicture}
\end{minipage}%
\hfill
\begin{minipage}{0.47\textwidth}
\centering
\begin{tikzpicture}[scale=0.55]
    \draw[thick, -Stealth] (0,0) -- (5,0) node[right, font=\tiny] {$X$};
    \draw[thick, -Stealth] (0,0) -- (0,5) node[above, font=\tiny] {$Y$};
    \foreach \p in {(0.5,4.5), (1,4.0), (1.5,3.6), (2,2.8), (2.5,2.5), (3,1.9), (3.5,1.4), (4,1.1), (4.5,0.5)} {
        \fill[blue!70] \p circle (2.5pt);
    }
    \draw[red, thick, dashed] (0.3,4.7) -- (4.7,0.3);
    \node[font=\scriptsize\bfseries] at (2.5,-0.8) {(b) $r_{xy} \approx -0.98$};
\end{tikzpicture}
\end{minipage}

\vspace{0.4cm}
\begin{minipage}{0.47\textwidth}
\centering
\begin{tikzpicture}[scale=0.55]
    \draw[thick, -Stealth] (0,0) -- (5,0) node[right, font=\tiny] {$X$};
    \draw[thick, -Stealth] (0,0) -- (0,5) node[above, font=\tiny] {$Y$};
    \foreach \p in {(1.2,3.5), (2.5,4.2), (3.8,3.2), (1.5,1.8), (2.8,2.2), (3.5,1.5), (2.0,2.8), (4.2,2.5), (1.0,2.5)} {
        \fill[blue!70] \p circle (2.5pt);
    }
    \node[font=\scriptsize\bfseries] at (2.5,-0.8) {(c) $r_{xy} \approx 0.0$ (Assenza correlazione)};
\end{tikzpicture}
\end{minipage}%
\hfill
\begin{minipage}{0.47\textwidth}
\centering
\begin{tikzpicture}[scale=0.55]
    \draw[thick, -Stealth] (0,0) -- (5,0) node[right, font=\tiny] {$X$};
    \draw[thick, -Stealth] (0,0) -- (0,5) node[above, font=\tiny] {$Y$};
    \foreach \p in {(0.5,4.0), (1.0,2.25), (1.5,1.0), (2.0,0.25), (2.5,0.05), (3.0,0.25), (3.5,1.0), (4.0,2.25), (4.5,4.0)} {
        \fill[red!80!black] \p circle (2.5pt);
    }
    \draw[red!80!black, thick] (0.5,4.0) parabola bend (2.5,0.05) (4.5,4.0);
    \node[font=\scriptsize\bfseries, align=center] at (2.5,-0.8) {(d) $r_{xy} \approx 0.0$ (Relazione quadratica!)};
\end{tikzpicture}
\end{minipage}
\caption{Interpretazione visiva del coefficiente di correlazione di Pearson: (a) correlazione positiva; (b) correlazione negativa; (c) assenza di relazione lineare; (d) dipendenza deterministica non lineare sfuggita a Pearson.}
\label{fig:pearson-cases}
\end{figure}

\subsection{Associazione Statistica vs Causalit\`a e Fattori di Confondimento}
Un principio fondamentale nella data science \`e la distinzione epistemologica tra associazione e causazione:
\begin{equation}
\textbf{Correlazione} \;\neq\; \textbf{Causalit\`a}
\end{equation}
Si consideri un dataset accademico in cui il numero di ore settimanali dedicate allo studio e i voti d'esame mostrano un'elevata correlazione positiva ($r_{xy} = +0.98$). Gli studenti con meno ore di studio tendono a trovarsi sotto la media, mentre chi studia pi\`u ore si colloca sopra la media. Tale dato riflette una fortissima associazione statistica, ma \textbf{non autorizza a concludere che 2 ore aggiuntive di studio causino direttamente e determiniscamente un incremento prefissato del voto}. Molteplici altre variabili (preparazione pregressa, frequenza alle lezioni, attenzione in aula) possono influenzare congiuntamente entrambi gli aspetti.

La correlazione lineare:
\begin{enumerate}
    \item Non identifica la direzione della causalit\`a ($X \to Y$ oppure $Y \to X$).
    \item Pu\`o essere generata da un \textbf{fattore di confondimento} (\textit{confounding variable}).
\end{enumerate}

\begin{example}[La Variabile Confondente: Gelati e Annegamenti]
Nelle statistiche estive si osserva una correlazione positiva molto forte tra il volume di vendite di gelati e il numero di incidenti da annegamento. Chiaramente mangiare gelati non provoca annegamenti. Esiste una terza variabile latente non inclusa nella coppia, la \textbf{temperatura esterna elevata}, che aumenta simultaneamente sia la domanda di gelati sia il numero di bagnanti in acqua, inducendo una correlazione spuria tra le due grandezze.
\end{example}

\subsection{Correlazione di Rango di Spearman per Relazioni Monotone}
Nei contesti in cui la dipendenza tra due grandezze non \`e strettamente lineare, oppure in presenza di outlier marcati che distorcono la media e la deviazione standard, si ricorre al \textbf{coefficiente di correlazione di rango di Spearman} ($\rho_s$).

Spearman non confronta direttamente i valori numerici grezzi, ma li sostituisce con la loro posizione ordinale o \textbf{rango} (\textit{rank}) nell'ordinamento crescente del campione:

\begin{definition}[Coefficiente di Correlazione di Rango di Spearman]
Siano $\operatorname{rank}(x_i)$ e $\operatorname{rank}(y_i)$ i ranghi di $x_i$ e $y_i$. Se non vi sono ex aequo (\textit{ties}), il coefficiente \`e formulato come:
\begin{equation}
\rho_s = 1 - \frac{6 \sum_{i=1}^n d_i^2}{n(n^2 - 1)}, \quad \text{con } d_i = \operatorname{rank}(x_i) - \operatorname{rank}(y_i)
\end{equation}
In presenza di ranghi parimerito, $\rho_s$ equivale al coefficiente di Pearson calcolato direttamente sulle sequenze dei ranghi.
\end{definition}

Propriet\`a essenziali:
\begin{itemize}
    \item $\rho_s = +1$ quando l'ordinamento delle due variabili coincide perfettamente (relazione \textbf{monotona crescente}).
    \item $\rho_s = -1$ quando l'ordinamento \`e esattamente invertito (relazione \textbf{monotona decrescente}).
    \item Misura se due grandezze tendono a variare nella stessa direzione o in direzione opposta, senza vincoli di linearit\`a (ad esempio per $Y = X^3$ o $Y = \log(X)$ si ha $\rho_s = 1$).
    \item Gode di elevata robustezza rispetto a valori estremi o outlier isolati.
\end{itemize}

\subsection{Matrice di Correlazione e Selezione delle Feature}
In dataset con molteplici attributi continui, si calcolano tutte le correlazioni bivariate a coppie organizzandole in una \textbf{matrice di correlazione} simmetrica $p \times p$.
Nell'analisi del dataset Iris (Figura~\ref{fig:iris-correlation-heatmap}), ad esempio, la matrice evidenzia una fortissima correlazione sia lineare che monotonica tra \textit{Petal Length} e \textit{Petal Width} ($r = 0.96$, $\rho_s = 0.94$), e tra \textit{Sepal Length} e \textit{Petal Length} ($r = 0.87$, $\rho_s = 0.88$), segnalando una rilevante ridondanza informativa che suggerisce la possibilit\`a di fondere o selezionare solo uno dei due attributi per l'addestramento di classificatori e modelli predittivi.

\begin{figure}[htbp]
\centering
\includegraphics[width=0.92\textwidth]{iris_correlation_heatmaps.pdf}
\caption{Matrici di correlazione bivariata calcolate sul dataset Iris: a sinistra la correlazione lineare di Pearson ($r$), a destra la correlazione di rango di Spearman ($\rho_s$). Si evidenzia visivamente l'intensa associazione positiva tra le dimensioni dei petali e l'associazione moderatamente negativa tra la larghezza del sepalo e le restanti feature.}
\label{fig:iris-correlation-heatmap}
\end{figure}

\section{Gestione di Outlier e Valori Mancanti}

\subsection{Outlier: Rumore vs Obiettivo Analitico}
Un \textbf{outlier} (anomalia o valore aberrante) \`e un'istanza con caratteristiche che deviano in modo marcato dalla maggior parte delle osservazioni del campione:
\begin{itemize}
    \item \textbf{Outlier come rumore (\textit{noise}):} Generati da errori materiali di digitazione, malfunzionamenti transitori di sensori fisici o problemi nella conversione dei formati. Tali osservazioni falsano drasticamente medie, varianze e modelli di regressione, e devono essere identificate e rimosse (o corrette) nel pre-processing.
    \item \textbf{Outlier come obiettivo analitico (\textit{target}):} Situazioni insolite o eccezionali reali, che rappresentano il primario scopo applicativo dell'analisi. Casi tipici comprendono la rilevazione di \textbf{frodi con carte di credito}, l'individuazione di \textbf{intrusioni informatiche di rete} (\textit{intrusion detection}) e la diagnosi precoce di \textbf{patologie mediche rare}.
\end{itemize}

Anche qualora gli outlier derivino da misurazioni corrette di eventi rari, pu\`o risultare opportuno escluderli temporaneamente dal calcolo di parametri aggregati qualora l'obiettivo sia caratterizzare il comportamento della popolazione tipica.

\subsection{Metodologie di Rilevamento degli Outlier}
L'identificazione delle anomalie si articola su differenti livelli di dimensionalit\`a:
\begin{itemize}
    \item \textbf{Attributo Singolo (Univariato):}
    \begin{itemize}
        \item \textit{Attributi categorici:} Un outlier corrisponde a un valore che si presenta con una frequenza sensibilmente inferiore rispetto a tutte le altre modalit\`a del dominio.
        \item \textit{Attributi numerici:} Ispezione visiva e analitica mediante Box Plot (valori esterni all'intervallo $[Q_1 - 1.5 \times \text{IQR}, \; Q_3 + 1.5 \times \text{IQR}]$) e calcolo dello $z$-score standardizzato ($z = \frac{x - \mu}{\sigma}$).
    \end{itemize}
    \item \textbf{Attributi Multipli (Multivariato):}
    \begin{itemize}
        \item \textit{Scatter Plot 2D e 3D:} Permettono di cogliere outlier sia globali (estremi rispetto all'intera nuvola di punti) sia locali (anomali rispetto al cluster o alla classe specifica a cui appartengono, come illustrato in Figura~\ref{fig:scatterplot-outlier}).
        \item \textit{Grafici PCA (\textit{Principal Component Analysis}):} Proiettano lo spazio multidimensionale sui piani di massima varianza per individuare punti isolati.
        \item \textit{Algoritmi di Clustering:} Negli algoritmi basati sulla densit\`a o su distanze (come DBSCAN o k-means con soglie di distanza), gli outlier emergono naturalmente come i punti che non possono essere assegnati ad alcun cluster principale o che rimangono completamente isolati.
    \end{itemize}
\end{itemize}

\subsection{Valori Mancanti (Missing Values) e Camuffamento}
L'assenza di valori per singoli attributi in determinate istanze pu\`o derivare da diverse cause:
\begin{itemize}
    \item \textbf{Informazione non raccolta o rifiutata:} Rifiuto esplicito dell'intervistato di condividere dati riservati (ad esempio reddito, peso corporeo o et\`a).
    \item \textbf{Non applicabilit\`a logica dell'attributo:} L'attributo non possiede significato logico per una specifica classe di istanze (ad esempio il reddito annuo per bambini o minori in et\`a scolare).
    \item \textbf{Malfunzionamento o blocco di sensori:} Interruzioni momentanee di connettivit\`a o dispositivi fisici danneggiati.
\end{itemize}

\begin{alertblock}{Insidia del Camuffamento: Zeri e Valori di Default Fittizi}
Spesso un valore mancante non viene salvato con codifica esplicita (\texttt{NULL} o \texttt{NaN}), bens\`i registrato con valori convenzionali o di default (ad esempio \texttt{0}, \texttt{-1} o \texttt{999}). Un sensore bloccato, ad esempio, invia ripetutamente il valore \texttt{0}.
Se non intercettati nella fase di Data Understanding, tali numeri vengono trattati acriticamente dai software come misurazioni reali, distorcendo gravemente medie, deviazioni standard e modelli di classificazione.
\end{alertblock}

\section{Checklist Operativa di Data Understanding}
Al termine della fase di esplorazione, l'analista deve aver completato sistematicamente le seguenti verifiche:
\begin{enumerate}
    \item \textbf{Valutazione della Qualit\`a dei Dati:} Verificare la presenza di errori di accuratezza sintattica (valori fuori dominio) e semantica, grado di completezza, tempestivit\`a, duplicati e grado di sbilanciamento delle classi.
    \item \textbf{Individuazione degli Outlier:} Rilevare anomalie univariati e multivariati mediante tecniche di visualizzazione (box plot, scatter plot, PCA, coordinate parallele) e decidere la politica di trattamento (rumore da rimuovere vs target primario di indagine).
    \item \textbf{Rilevamento ed Esame dei Valori Mancanti:} Identificare i campi nulli e smascherare valori mancanti nascosti o camuffati sotto forma di zeri fittizi o codifiche di default.
    \item \textbf{Scoperta e Conferma di Correlazioni e Dipendenze:} Calcolare matrici di correlazione lineare (Pearson) e di rango (Spearman) per confermare relazioni attese, rilevare dipendenze inattese e scongiurare la multicollinearit\`a e le ridondanze tra feature.
    \item \textbf{Verifica di Ipotesi e Assunzioni Specifiche del Dominio:} Controllare assunzioni applicative e distributive (ad esempio verificare se un attributo segue una distribuzione approssimativamente normale per l'uso della regola di Sturges o di stimatori parametrici).
    \item \textbf{Confronto delle Statistiche con il Comportamento Atteso:} Verificare che medie, mediane, asimmetrie e dispersione riflettano coerentemente la conoscenza di dominio e i pattern attesi nel contesto applicativo.
\end{enumerate}
